Is position bias a property of small judges? We ran the larger Qwen3-VL-8B-Instruct on the same 300 prompts and three orders. To fit it in memory we used 8-bit weights, so we also reran the 4B judge with the same quantization as a control. The control agrees with the main 4B run on 94.9\% of calls and reproduces its first-slot share (48.2\%) and gain (0.044), so neither quantization nor the MLX runtime explains the differences below (Figure~\ref{fig:scale}, Appendix Table~\ref{tab:scale}).

At 8B the position bias largely disappears: the first slot is chosen in 27.3\% of calls, close to the uniform 27.8\% ($p=0.06$), although the choice still changes across orders on 43.0\% of prompts. The 8B judge is also much better: its gain over random is 0.097 [0.077, 0.117], above the CLIP baseline (0.069; paired difference 0.028 [0.003, 0.054]), and its stereotype difference is no longer detectable (0.016 [$-$0.004, 0.036]).

The value of the unanimity gate changes with it. For the 4B judge, the prompts unanimity rejects fall below random; for the 8B judge they stay above random (0.068 [0.036, 0.099]), so the gate mostly discards good decisions. Agreement between the 8B and 4B judges shows the same pattern: on the 126 prompts where they disagree, the 8B choice has gain 0.098 and the 4B choice $-$0.039. A gate is useful only while it filters the failure mode of the judge in front of it, and must be re-audited when the judge changes.
